\section{Related Work}
\label{section:related_work}

We discuss connections to related work in queuing theory, reinforcement learning, and gradient estimation in machine learning and operations research.

\paragraph{Scheduling in Queuing Networks}

Scheduling is a long-studied control task in the queuing literature for managing queues with multiple classes of jobs ~\cite{harrison1989scheduling, meyn2008control}. Standard policies developed in the literature include static priority policies such as the $c\mu$-rule~\cite{cox1961queues}, threshold policies~\cite{rosberg1982optimal}, policies derived from fluid approximations~\cite{avram1994optimal, chen1993dynamic, meyn2001sequencing}, including discrete review policies~\cite{harrison1996bigstep, maglaras2000discrete}, policies that have good stability properties such as MaxWeight ~\cite{stolyar2004maxweight} and MaxPressure ~\cite{dai2005maximum}. Many of these policies satisfy desirable properties such as throughput optimality~\cite{tassiulas1990stability, armony2003queueing}, or cost minimization~\cite{cox1961queues, mandelbaum2004scheduling} for certain networks and/or in certain asymptotic regimes. In our work, we aim to leverage some of the theoretical insights developed in this literature to design reinforcement learning algorithms that can learn faster and with less data than model-free RL alternatives. We also use some of the standard policies as benchmark policies when validating the performance of our $\pathwise$ policy gradient algorithm.

\paragraph{Reinforcement Learning in Queueing Network Control}
Our research connects with the literature on developing reinforcement learning algorithms for queuing network control problems~\citep{moallemi2008approximate, shah2020stable, qu2020scalable, dai2022queueing, liu2022rl, wei2024sample, pavse2024learning}. These works apply standard model-free RL techniques (e.g. $Q$-learning, $\ppo$, value iteration, etc.) but introduce novel modifications to address the unique challenges in queuing network control problems. Our work differs in that we propose an entirely new methodology for learning from the environment based on differentiable discrete event simulation, which is distinct from all model-free RL methods. The works~\cite{shah2020stable,liu2022rl, dai2022queueing, pavse2024learning} observe that RL algorithms tend to be unstable and propose fixes to address this, such as introducing a Lyapunov function into the rewards, or behavior cloning of a stable policy for initialization. In our work, we propose a simple modification to the policy network architecture, denoted as the \emph{work-conserving softmax} as it is designed to ensure work-conservation. We find empirically that \emph{work-conserving softmax} ensures stability with even randomly initialized neural network policies. In our empirical experiments, we primarily compare our methodology with the $\ppo$ algorithm developed in~\cite{dai2022queueing}. In particular, we construct a $\ppo$ baseline with the same hyper-parameters, neural network architecture, and variance reduction techniques as in~\cite{dai2022queueing}, although with our policy architecture modification that improves stability. 
% It is also worth noting that these works typically benchmark their policies in environments with exponentially-distributed inter-arrival and service times, in order to use the discrete MDP representation of the queuing network. Our framework, based in discrete-event simulation, is capable of modeling general inputs and we benchmark our policies in environments with higher-variance noise.

\paragraph{Differentiable Simulation in RL and Operations Research}

While differentiable simulation is a well-studied paradigm for control problems in physics and robotics~\cite{heiden2021neuralsim, hu2019difftaichi, suh2022differentiable, howell2022dojo, schoenholz2020jax}, it has only recently been explored for large-scale operations research problems. For instance, \cite{madeka2022deep,alvo2023neural} study inventory control problems and train a neural network using direct back-propagation of the cost, as sample paths of the inventory levels are continuous and differentiable in the actions. In our work, we study control problems for queuing networks, which are discrete and non-differentiable, preventing the direct application of such methods. To address this, we develop a novel framework for computing pathwise derivatives for these non-differentiable systems, which proves highly effective for training control policies. Another line of work, including \cite{andelfinger2021differentiable, andelfinger2023towards}, proposes differentiable agent-based simulators based on differentiable relaxations. While these relaxations have shown strong performance in optimization tasks, they also introduce unpredictable discrepancies with the original dynamics. We introduce tailored differentiable relaxations in the back-propagation process only, ensuring that the forward simulation remains true to the original dynamics. 

\paragraph{Gradient Estimation in Machine Learning}

Gradient estimation~\cite{mohamed2020monte} is an important sub-field of the machine learning literature, with applications in probabilistic modeling~\cite{kingma2013auto, jang2016categorical} and reinforcement learning~\cite{williams1992simple, sutton1999policy}. There are two standard strategies for computing stochastic gradients~\cite{mohamed2020monte}. The first is the score-function estimator or $\reinforce$~\cite{williams1992simple,sutton1999policy}, which only requires the ability to compute the gradient of log-likelihood but can have high variance~\cite{greensmith2004variance}. Another strategy is the reparameterization trick~\cite{kingma2013auto}, which involves decomposing the random variable into the stochasticity and the parameter of interest, and then taking a pathwise derivative under the realization of the stochasticity. Gradient estimators based on the reparameterization trick can have much smaller variance~\cite{mohamed2020monte}, but can only be applied in special cases (e.g. Gaussian random variables) that enable this decomposition. Our methodology makes a novel observation that for queuing networks, the structure of discrete-event dynamical systems gives rise to the reparameterization trick. Nevertheless, the function of interest is non-differentiable, so standard methods cannot be applied. As a result, our framework also connects with the literature on gradient estimation for discrete random variables~\cite{jang2016categorical, maddison2016concrete,  bengio2013estimating, tucker2017rebar}. In particular, to properly smooth the non-differentiability of the event selection mechanism, we employ the straight-through trick~\cite{bengio2013estimating}, which has been previously used in applications such as discrete representation learning~\cite{van2017neural}. Our work involves a novel application of this technique for discrete-event systems, and we find that this is crucial for reducing bias when smoothing over long time horizons.


\paragraph{Gradient Estimation in Operations Research}
There is extensive literature on gradient estimation for stochastic systems~\cite{glasserman1990gradient, glasserman1992derivative, glynn1987likelilood, cao1985convergence, fu2012conditional}, some with direct application to queuing optimization ~\cite{l1994stochastic, glynn1990likelihood, fu2012conditional}.
Infinitesimal Perturbation Analysis (IPA)~\cite{glasserman1990gradient, ho1983infinitesimal, johnson1989infinitesimal} is a standard framework for constructing pathwise gradient estimators, which takes derivatives through stochastic recursions that represent the dynamics of the system. While IPA has been applied successfully to some specific queuing networks and discrete-event environments more broadly~\cite{suri1987infinitesimal}, standard IPA techniques cannot be applied to general queuing networks control problems, as has been observed in~\cite{cao1987first}. There has been much research on outlining sufficient conditions under which IPA is valid, such as the commuting condition in~\cite{glasserman1990gradient, glasserman1992derivative} or the perturbation conditions in~\cite{cao1985convergence}, but these conditions do not hold in general. Several extensions to IPA have been proposed, but these alternatives require knowing the exact characteristics of the sampling distributions and bespoke analysis of event paths~\cite{fu2012conditional, fu1994smoothed}. Generalized likelihood-ratio estimation~\cite{glynn1987likelilood} is another popular gradient estimation framework, which leverages an explicit Markovian formulation of state transitions to estimate parameter sensitivities. However, this requires knowledge of the distributions of stochastic inputs, and even with this knowledge, it may be difficult to characterize the exact Markov transition kernel of the system. Finally, finite differences~\cite{fu1997optimization} and finite perturbation analysis~\cite{ho1983infinitesimal, cao1987first} are powerful methods, particularly when aided with common random numbers~\cite{glynn1989optimization, glasserman1992some}, as it requires minimal knowledge about the system. However, it has been observed that performance can scale poorly with problem dimension~\cite{glasserman2004monte, glynn1989optimization}, and we also observe this in an admission control task (see Section~\ref{section:admission}).

Our contribution is proposing a novel, general-purpose framework for computing pathwise gradients through careful smoothing, which only requires samples of random input (e.g., interarrival times and service times) rather than knowledge of their distributions.
%and can be applied to any multi-class queuing problem under any (differentiable) scheduling policy. 
Given the negative results about the applicability of IPA for general queuing network control problems (e.g., general queuing network model and scheduling policies), we introduce bias through smoothing to achieve generality. It has been observed in~\cite{eckman2020biased} that biased IPA surrogates can be surprisingly effective in simulation optimization tasks such as ambulance base location selection. Our extensive empirical results confirm this observation and illustrate that while there is some bias, it is very small in practice, even over long time horizons ($>10^5$ steps).

% provides a novel and efficient way to compute/approximate the pathwise gradients through careful smoothing. It can be thought of as a biased IPA estimator with carefully controlled bias. It has been shown in~\cite{eckman2020biased} that despite non-differentiability issues that prevent the use of standard IPA estimators, biased IPA surrogates are surprisingly effective in simulation optimization tasks such as ambulance base location selection. In the following sections, we demonstrate through extensive numerical experiments that even though our proposed estimator is biased, the bias is very small in practice even over long time horizons ($>100,000$ steps). In addition, unlike many existing gradient estimators, our approach is general and scalable, which allows us to compare our approach to standard model-free gradient estimators like $\reinforce$ across a wide range of large-scale settings.

